\section{Introduction}\label{sec:intro}

\subsection{The Palais--Smale condition}\label{sec:palais-smale-cond}
A Palais--Smale sequence (or (PS)-sequence) for a continuously differentiable
energy function $\EL$ on a smooth Banach manifold $\ML$, is a sequence
$(x_k)_k\subset\ML$ such that $\sup_{k\in\N}\EL(x_k)<\infty$ and
$\lim_{k\to\infty}\|\D \EL_{x_k}\|_{(T_{x_k}\ML)^*}=0$, where 
$(T_x\ML)^*$ denotes the dual space of the tangent space $T_x\ML$ of 
the manifold $\ML$ at the point $x\in\ML$.

The energy $\EL$ is said to satisfy the \emph{Palais--Smale condition}
or (PS)-\emph{condition}, if every (PS)-sequence $(x_k)_k\subset\ML$
for $\EL$ contains a convergent subsequence, $x_{k_l}\to x_\infty\in\ML$
as $l\to\infty$.

The (PS)-condition has been central for the investigation of critical
points and of energy landscapes, with sophisticated tools such as
Mountain Pass theorems, 
Ljusternik-Schnirelman theory, or Morse-Bott results; 
see, e.g., the monographs \cite{rabinowitz_1986}, \cite{palais-terng_1988},
\cite{mawhin-willem_1989}, \cite{WillemMinimaxTheorems1996a},
\cite[Chapters II \& III]{struwe_2008a}, and
\cite{jabri_2003}. 

In \emph{geometric knot theory} one considers non-local energies that
penalize self-inter\-sec\-tions of geometric objects like curves or surfaces, 
so called \emph{self-avoidance energies} or \emph{knot energies}.
With these functionals one can treat variational problems in spaces of
embeddings under topological constraints, such as a prescribed knot type.
In this way, energy minimizing curves or Euclidean submanifolds have been
found in given ambient isotopy classes. In addition, bounds on the number of
isotopy types under given energy values were established, and classical knot
invariants such as the crossing number or the stick number can be controlled
in terms of knot energies; see the surveys 
\cite{ohara_2003,strzelecki-etal_2013a,
strzelecki-vdm_2013a,strzelecki-vdm_2014,blatt-reiter_2014,strzelecki-vdm_2018} 
and the many references 
therein. Moreover, critical points
were found by 
symmetric criticality, not only for self-avoidance energies   alone
\cite{cantarella-etal_2014b,gilsbach-vdm_2018,wings_2025}, 
but also for linear combinations of such
with the classic Euler-Bernoulli bending energy
\begin{equation}\label{eq:bending-energy}
        \BendingEnergy(\g) \ceq \frac{1}{2} \int_\g\kappa^2\d s
\end{equation}
for closed curves $\g:\R/\Z\to\R^3$, in search of \emph{elastic knots}
\cite{gerlach-etal_2017,gilsbach-etal_2023}. There are long-time existence
results for different kinds of  gradient flows for various knot energies, the 
$L^2$-gradient flow 
\cite{blatt_2012b,blatt_2018,blatt_2020a}, and Hilbert or Banach gradient flows
\cite{knappmann-etal_2023,matt-etal_2023}, as well as
$L^2$-gradient flows
for linear combinations with the bending energy \cref{eq:bending-energy}
in \cite{lin-schwetlick_2010,vonbrecht-blair_2017}. We should also point
out the substantial advances in the numerical treatment of such 
problems as in 
\cite{bartels-etal_2018,bartels-reiter_2020,bartels-reiter_2021,yu-etal_2021b,yu-etal_2021,sassen-etal_2024}.

However, a more sophisticated critical point theory for knot energies
has not been developed yet, partially because of the lack of the (PS)-condition
for these non-local energies. With this paper we want to close this gap
for a family of knot energies and for linear combinations of basically
all known differentiable knot energies
with the bending energy \cref{eq:bending-energy}.

\subsection{Three families of knot energies on curves}\label{sec:knot-energies}
O'Hara introduced and investigated
in \cite{ohara_1991a,ohara_1992a,ohara_1992b,ohara_1994}
\emph{repulsive potentials} of the form
\begin{equation}\label{eq:E-alpha-p}
E^{\alpha,p}\ceq\int_{\R/\Z}\int_{\R/\Z}\left(\frac1{|\g(u)-\g(v)|^\alpha}-\frac1{d_\g(s,t)^\alpha}\right)^p|\g'(u)||\g'(v)|\d u\d v,
\end{equation}
where $d_\g(u,v)$ denotes the intrinsic distance between the points $\g(u)$
and $\g(v)$ along the curve, regularizing the singularity of the first
summand. The most prominent example is the M\"obius energy
$E^{2,1}$, which  is not only scale-invariant but also invariant under
M\"obius transformations \cite[Theorem 2.1]{freedman-etal_1994}.

In contrast to O'Hara's energies \cref{eq:E-alpha-p} the \emph{(generalized)
tangent-point energies}
\begin{equation}\label{eq:TP}
 \TP^{(p,q)}(\gamma)= \int_{\R/\Z}\int_{-\frac{1}{2}}^\frac{1}{2}
\frac{\big|\P_{\gamma'(u)}^{\perp}(\gamma(u+w)-\gamma(u)) \big|^q}{|
\gamma(u+w)-\gamma(u)|^{p}}|\g'(u)||\g'(u+w)|\d w\d u
\end{equation}
have integrands that are finite for smooth embedded curves $\g \colon \R/\Z \to \R^n$
for a reasonable\footnote{The relevant ranges for the parameters of all
three energy families will be discussed below, in connection with
their respective energy spaces.}
parameter range for $p$ and $q$.
So, there is no need for a regularizing term. Here, $P_v^\perp$ means the
orthogonal projection onto the $(n-1)$-dimensional
linear subspace $v^\perp\subset\R^n$ for a vector $v\in\R^n\setminus\{0\}$.
The energies' name originates in the
geometric classic tangent-point energy $\TP_q\ceq2^q\TP^{(2q,q)}$ whose
integrand equals the $q$-th power of the inverse
tangent-point radius, i.e., the
radius of the unique circle through the points
$\g(u)$ and $\g(u+w)$ that is in addition tangent to the curve
$\g$ in $\g(u)$. 
This energy was suggested by Gonzalez and Maddocks in 
\cite[p. 4773]{gonzalez-maddocks_1999} as an alternative to self-repulsive
potentials, and was first investigated analytically in
\cite{strzelecki-vdm_2012}.

In addition to the tangent-point radius that defines $\TP_q$,
Gonzalez
and Maddocks suggested to also
consider the unique circle through three given curve points, which
leads to integral Menger curvature $\ML_p\ceq2^p\intM^{(p,p)}$
as a specific member of the family of
\emph{(generalized) integral Menger curvatures}
\begin{equation}\label{eq:intM}
\Mpq=\iiintRZ\frac{\left|\gamma'(u) \right|\left|\gamma'(v) \right|\left|\gamma'(w) \right|}{R^{(p,q)}(\gamma(u),\gamma(v),\gamma(w))}\d u \d w \d v,
\end{equation}
where
\begin{equation}\label{eq:integrand-intM}
R^{(p,q)}(x,y,z)\ceq\frac{\left( \left|y-z \right|\left|y-x \right|\left|z-x \right| \right)^p}{\left|(y-x)\wedge(z-x) \right|^q}\, \text{ for }x,y,z\in\R^n;
\end{equation}
see \cite{strzelecki-etal_2010} for first analytical
results on $\ML_p$ and \cite{strzelecki-etal_2013a} for its
impact on knot theory. 

Crucial for our investigations is the knowledge of the
exact \emph{energy spaces} for these families of knot energies.
By this we mean specific Banach spaces $\CL$ such that any arclength
parametrized closed curve has finite energy if and only if
it is embedded and contained in $\CL$.
For each of these energies the underlying energy space\footnote{For $\TP^{(p,q)}$
this was shown in \cite[Theorem 1.1]{blatt-reiter_2015a}, for $\intM^{(p,q)}$
see \cite[Theorem 1]{blatt-reiter_2015b}, and for $E^{\alpha,p}$ we 
refer to \cite[Theorem 1.1]{blatt_2012a}.}
 is a fractional
Sobolev space $W^{1+s,\rho}(\R/\Z,\R^n)$, where
\begin{equation}\label{eq:individual-energy-spaces}
\textstyle
(s,\rho)\ceq\begin{cases}
\big(\frac{\alpha p-1}{2p},2p\big) & \Fo \FL=E^{\alpha,p},\,
\alpha>0,p\ge 1,2<\alpha p < 2p+1.\\
\big(\frac{p-1-q}{q},q\big) & \textstyle
\Fo \FL=\TP^{(p,q)},\,p\in (q+2,2q+1),q>1,\\ 
\big(\frac{3p-2q-2}q,q\big) & \textstyle
\Fo \FL=\intM^{(p,q)},\,
p\in(\frac23 q+1,q+\frac23), q>1.
\end{cases}
\end{equation}
Note that in all these cases, $s\in (0,1)$.
If $\rho=2$ these fractional Sobolev spaces are Hilbert spaces, and we write
$W^{1+s,2}=H^{1+s}$, where we also admit
the case $s=1$ to include the bending energy \cref{eq:bending-energy},
whose energy space equals $H^2(\R/\Z,\R^n)$. In fact, in the present
paper we focus on the case when the energy space of the
leading order energy is a Hilbert space, allowing, e.g., for
the tangent-point energies $\TP^{(p,2)}$ for $p\in (4,5)$,
and integral Menger curvature $\intM^{(p,2)}$ for $p\in (\frac73,\frac83)$,
as well as total energies of the form $\TL \ceq \BendingEnergy + \theta \, \FL$ for some
$\theta>0$, where $\FL$ is any of the energies in
\cref{eq:individual-energy-spaces}.

Since the knot energies and the bending energy are invariant under the
action of the non-compact group of reparametrizations of the domain $\R/\Z$,
we restrict 
them to the set of injective, arclength 
parametrized curves,
\begin{equation}\label{eq:ALs}
        \AL^{1+s}
        \ceq
        \braces[\big]{
                \g\in H^{1+s}(\R/\Z,\R^n):\g_{[0,1)}\,\,\textnormal{injective}, \,|\g'|=1
        }
        ,
\end{equation}
and show in \Cref{thm:arclength-manifold} that $\AL^{1+s}$ is
a smooth submanifold of  $H^{1+s}(\R/\Z,\R^n)$ for any $s\in (\frac{1}{2},1]$.
It contains another smooth submanifold 
\begin{equation}\label{eq:ALsz}
        \textstyle
        \AL^{1+s}_\bz
        \ceq
        \braces[\big]{\g\in\AL^{1+s}:\g(0)=\bz},\quad s\in (\frac{1}{2},1],
\end{equation}
where $\bz\in\R^n$ is an arbitrary fixed point;
see \Cref{thm:arclength-manifold-point}. 
When restricting the energies
to these submanifolds we do not lose any topological information, since
we prove in \Cref{thm:strong-deformation-ALs} that the open subset of all
regular knots in $H^{1+s}(\R/\Z,\R^n)$ strongly deformation retracts
onto the submanifolds $\AL^{1+s}$ and $\AL^{1+s}_\bz$ 
for every $s\in (\frac{1}{2},1]$
and $\bz\in\R^n$. 
In addition, if two knots $\g_0$ and $\g_1$ lie in
the same connected component of $\AL^{1+s}$ then they are ambiently isotopic,
i.e., in the same knot class,
as shown in \Cref{thm:components-knot-classes}.


\subsection{Main results}
We start with the Palais--Smale condition for linear combinations of the
bending energy \cref{eq:bending-energy} with any energy contained in
one of the
three families \cref{eq:E-alpha-p}, \cref{eq:TP}, or \cref{eq:intM}. But we
have to shrink slightly\footnote{This does \emph{not} change the 
respective parameter range
for $E^{\alpha,1}$, or for $\TP^{(p,q)}$ and $\intM^{(p,q)}$ if $q\in (1,2]$.}
the range of parameters by requiring in addition to
\cref{eq:individual-energy-spaces} that
\begin{equation}\label{eq:shrink-range}
\textstyle
\alpha p<p+2\Fo E^{\alpha,p},\,\, p<\frac{3}{2} q +2\Fo
\TP^{(p,q)},\AND p<\frac{5}{6} q+1\Fo \intM^{(p,q)}
\end{equation}
to allow for the compact embedding $H^2\hookrightarrow W^{1+s,\rho}$ for 
each knot energy. 
\begin{theorem}[Bending energy plus knot energy above scale-invariance]
\label[theorem]{thm:bending-knot-energy}
Let $\theta>0$,  $\bz\in\R^n$, and $\FL\in\{E^{\alpha,p},\TP^{(p,q)},
\intM^{(p,q)}\}$, where the respective combinations of
parameters $\alpha$, $p$, and $q$ are chosen
as in \cref{eq:individual-energy-spaces} and \cref{eq:shrink-range}.
Then the total energy $\TL \ceq \BendingEnergy + \theta \, \FL$ restricted to
the arclength submanifold $\AL^2_\bz\subset H^2(\R/\Z,\R^n)$ 
defined in \cref{eq:ALsz}
satisfies the {\rm (PS)}-condition.
\end{theorem}
Notice that \cref{eq:individual-energy-spaces} excludes the respective
scale-invariant energies, whose energy spaces do \emph{not} embed compactly
into $C^1(\R/\Z,\R^n)$. For example, the energy space of the
M\"obius energy $E^{2,1}$ is the Hilbert space
$H^{\frac{3}{2}}(\R/\Z,\R^n)$. To include this energy as part of
the total energy, we
use in our proof
a slightly smaller Banach space that \emph{does} embed
into $C^1$ (similarly as in 
\cite{reiter-schumacher_2021} and \cite{matt-etal_2023}).
\begin{theorem}[Bending energy plus M\"obius energy]
\label[theorem]{thm:bending-moebius-energy}
The total energy $\TL \ceq \BendingEnergy+ \theta \, E^{2,1}$ restricted
to $\AL^2_\bz$ satisfies the {\rm (PS)}-condition for every $\theta>0$ and
$\bz\in\R^n$.
\end{theorem}
We strongly believe that the statement of \Cref{thm:bending-knot-energy}
carries over to all scale-invariant energies, i.e., to $E^{\alpha,p}$
for $\alpha p=2$, to $\TP^{(p,q)}$ for $p=q+2$, and to $\intM^{(p,q)}$
for $p=\frac23 q +1$, as well.

The case of a pure knot energy  is harder to handle, and we focus
here on the tangent-point energies $\TP^{(p,2)}$ for $p\in (4,5)$ whose
energy space $H^{1+\frac{p-3}2}(\R/\Z,\R^n)$ is a Hilbert space.
\begin{theorem}[Tangent-point energies $\TP^{(p,2)}$]
\label{thm:PS-TP}
Let $\bz\in\R^n$ and $p\in (4,5)$ be fixed and set
\begin{equation}\label{eq:s-TP}
        \textstyle
        s(p)\ceq\frac{p-3}{2}.
\end{equation}
Then the tangent-point energy $\TP^{(p,2)}$ restricted to
the arclength submanifold $\AL^{1+s(p)}_\bz\subset H^{1+s(p)}(\R/\Z,\R^n)$
satisfies the {\rm (PS)}-condition.
\end{theorem}
This result carries over to O'Hara's energies $E^{\alpha,1}$ for 
$\alpha\in (2,3)$ and to the integral Menger curvature functionals
$\intM^{(p,2)}$ for $p\in (\frac73,\frac83)$,   where
one has to replace \cref{eq:s-TP} by the respective differentiability
index $s=\frac{\alpha-1}2$ or 
$s=\frac{3}{2} p -3$ of the corresponding energy space;
see \cite{freches_2025}. At this point, it seems feasible to generalize our methods to the non-Hilbert case $ \TPpq $, with suitable restrictions on $ p $ and $ q $. This is currently work in progress (see \cite{frechesetal_2026}). 

As an immediate application of the (PS)-condition we can
prove the existence of minimizing (PS)-sequences and of
minimizing knots for the total energies
$\TL = \BendingEnergy + \theta \, \FL$ of \Cref{thm:bending-knot-energy}
in any connected component of the arclength
submanifold.

\begin{corollary}[Minimizing knots for total energies]
\label{cor:minimizing-knots-total-energy}
Fix a point \,$\bz\in\R^n$ and some $\theta>0$, and let $\TL = \BendingEnergy + \theta \, \FL$
be any of the total energies considered 
in \cref{thm:bending-knot-energy,thm:bending-moebius-energy}. 
Let $\KL$ be a connected component of the arclength submanifold $\AL^2_\bz$.
Then there are a minimizing {\rm (PS)}-sequence and
a minimizing knot for $\TL|_{\KL}$.
The same statements hold true for the whole arclength submanifold, i.e., for $\KL = \AL^2_\bz$.
\end{corollary}
This result contains the existence results in
\cite[Theorem 3.1]{vdm_1999} and \cite[Theorem 1.1]{gilsbach-etal_2023}\footnote{The additional symmetry constraint considered in \cite[Theorem 1.1]{gilsbach-etal_2023} is of lower order and can be omitted to obtain
global minimizers. This was carried
out  in \cite[Theorem 2.1]{gerlach-etal_2017} but there
for the non-smooth ropelength
functional, which  can be viewed 
as the $L^\infty$-limit of $\TP_q^{1/q}$, or of
$\ML_p^{1/p}$ as $q$ or $p$ tends to infinity.}, there
for specific
choices of knot energies
in pursuit of \emph{elastic knots} in any given knot class.  
As mentioned before,
connected components of $\AL^2_\bz$ correspond to
different knot classes as we show in \Cref{thm:components-knot-classes}.

For the pure  tangent-point energies $\TP^{(p,2)}$,
which are even
of class $C^{1,1}_\loc$ on 
 the underlying arclength submanifold of the larger Hilbert space $H^{1+s(p)}
 (\R/\Z,\R^n)$,
we obtain the analogous statement.


\begin{corollary}[Minimizing knots for $\TP^{(p,2)}$]
\label{cor:minimizing-knots-TP}
Let $\bz\in\R^n$, $p\in (4,5)$, and $s(p)$ as in \cref{eq:s-TP}.
Let $\KL$ be a connected component of the arclength submanifold $\AL^{1+s(p)}_\bz$.
Then there are a minimizing {\rm (PS)}-sequence and a minimizing knot for $\TP^{(p,2)}|_{\KL}$.
This is also true for $\KL = \AL^{1+s(p)}_\bz$.
\end{corollary}
The existence of minimizing knots in prescribed knot classes
under a fixed-length constraint
has been proven in \cite[Theorem 1.3]{blatt-reiter_2015a} for the 
more general energy family
$\TP^{(p,q)}$ with parameters $p,q$ as in \cref{eq:individual-energy-spaces}.

If the energies are of class $C^{1,1}_\loc$, then one can apply
Hilbert gradient-flow techniques on the arclength submanifold
to find critical points. This was carried out by Okabe and Schrader for
the pure bending energy $\BendingEnergy$ in \cite{okabe-schrader_2023}
 which inspired the following result.
Since at this point,  the necessary 
Lipschitz-estimates for the first
variation are not available
for \emph{every} knot energy considered above,
we restrict the following result
to the subfamilies of energy functionals $\FL$ with 
\begin{align}\label{eq:subfamilies}
\FL=E^{2,1},&\OR 
\textstyle
\FL\in 
\big\{\intM^{(p,2)} : p\in (\frac{7}{3} , \frac{8}{3}) \big\}
\cup
\big\{\intM^{(p,p)} : p\in (3,6)\big\},
\notag\\
&\OR\FL\in\big\{\TP^{(p,q)} : \textnormal{$p,q$ as in \cref{eq:individual-energy-spaces},
\cref{eq:shrink-range}, and $q\ge 2$}\big\}. 
\end{align}
\begin{theorem}[Gradient flow for total energies]
\label{thm:gradient-flow-total-energy}
        Fix $\theta >0$ and $\bz\in\R^n$, set $\ML\ceq\AL^2_\bz$,
        and consider the restriction
        $\widetilde{\TL}\ceq\TL|_{\ML}$ of the total energy
        $\TL \ceq \BendingEnergy + \theta \, \FL$ to the arclength submanifold $\ML$
        for any choice of $\FL$ in \cref{eq:subfamilies}.
        Then there is a function $\phi:[0,\infty)\times\ML\to\ML$ such
        that $\phi(\cdot,\g_0)\in C^1([0,\infty),\ML)$ for every knot
        $\g_0\in\ML$, that solves the Cauchy-Problem 
        \begin{equation}
                \label{eq:gradient-flow-total-energy}
                \begin{cases}
                        \frac{\d}{\d t}\phi(t,\g_0)& =-\nabla
                        \widetilde{\TL}(\phi(t,
                        \g_0))   
                                \Foa t\in [0,\infty),\\
                        \phi(0,\g_0)&=\g_0.  
                \end{cases}
        \end{equation}
        The function $t\mapsto\widetilde{\TL}(\phi(t,\g_0))$ 
        is strictly decreasing on $[0,\infty)$ unless
        $\g_0$ is 
        a critical knot for $\widetilde{\TL}$.
        For every $t\in [0,\infty)$ the knot $\phi(t,\g_0)$ is ambiently isotopic to $\g_0$, and as
        $t\to\infty$ the solution $\phi(t,\g_0)$ subconverges
        to a critical point $\g_\infty$ of  
        $\widetilde{\TL}$, which is ambiently isotopic
        to $\g_0$ as well. 
\end{theorem}  

This Hilbert gradient flow result is new, and it 
complements the long-time existence theorems
for the $L^2$-gradient flow of $\BendingEnergy+E^{2,1}$ in 
\cite[Theorem 1]{lin-schwetlick_2010},
as well as the arclength-constrained $L^2$-gradient flow in
\cite[Theorem 4.13]{vonbrecht-blair_2017} for total energies with a certain
family
of self-repulsive
potentials containing $E^{2,1}$.

Note, that for geometric energies, one usually investigates a natural gradient flow with respect to some geometric metric, often denoted $ L^2(\dd s) $, or $ H^2(\dd s) $. However, since we restrict the $ H^2 $-flow \cref{eq:gradient-flow-total-energy} to the arclength manifold, it coincides with the projected (i.e. intrinsic) $ H^2(\d s) $-gradient flow.

The subfamilies of energies in \cref{eq:subfamilies} contain 
the tangent-point energies $\TP^{(p,2)}$ for $p\in (4,5)$. For these
pure knot energies we can
also establish long-time existence for the Hilbert gradient flow
and its subconvergence
to critical points.
\begin{theorem}[Gradient flow for $\TP^{(p,2)}$]
\label{thm:gradient-flow-TP}
For fixed $\bz \in \R^n$, $p\in (4,5)$, and $s(p)$ as in \cref{eq:s-TP} set 
$\ML\ceq\AL^{1+s(p)}_\bz\subset H^{1+s(p)}(\R/\Z,\R^n)$. 
Then for  any given initial knot $\g_0\in\ML$
we have long-time existence for the gradient
flow \cref{eq:gradient-flow-total-energy} on $\ML$ 
within the same knot class $[\g_0]$ (with $\widetilde{\TL}$ 
replaced by the restriction $\widetilde{\TP}\ceq
\TP^{(p,2)}|_{\ML}$), along which $\widetilde{\TP}$ is strictly
decreasing unless the initial knot $\g_0$ is 
a $\widetilde{\TP}$-critical point. Moreover,
the flow subconverges to a $\widetilde{\TP}$-critical knot $\g_\infty$
of the  same
knot type as $\g_0$ as time
tends to infinity. 
\end{theorem}
This result carries over to the integral Menger curvatures
$\intM^{(p,2)}$ for $p\in (\frac{7}{3},\frac{8}{3})$ and to $ E^{\alpha,1} $ for $ \alpha \in (2,3) $, \cite{freches_2025}.
In contrast to the long-time existence result for a projected 
Hilbert-gradient
flow for $\intM^{(p,2)}$ in
\cite[Theorem 1.1]{knappmann-etal_2023} we do not need to project
the flow to control the parametrizations, and we do have subconvergence
of the flow
to a critical knot. Very recently, D\"ohrer in cooperation with the
first author proved a {\L}ojasiewicz--Simon inequality, which even leads
to full 
convergence of the flow for $\TP^{(p,2)}$ in \Cref{thm:gradient-flow-TP}; 
see \cite[Theorem 1.5]{doehrer-freches_2025}.
In addition, the gradient flows in \cref{thm:gradient-flow-total-energy,thm:gradient-flow-TP} preserve
symmetries, induced, for instance, by 
a group of isometries acting on the arclength submanifold; see 
\Cref{prop:symmetry-preserving-grad-flow} where we address such symmetries
in an abstract setting. 

In view of the existing regularity results\footnote{ 
\cite[Theorem 1.5]{blatt-reiter_2015a}, \cite[Theorem 4]{blatt-reiter_2015b}, \cite[Theorem 1.2]{blatt-reiter_2013}}
for critical points of knot energies
on the open subset of regular knots under a fixed-length
constraint 
one may ask whether also the more constrained arclength critical knots obtained in
\cref{cor:minimizing-knots-total-energy,cor:minimizing-knots-TP} and \cref{thm:gradient-flow-total-energy,thm:gradient-flow-TP} are smooth. 
We prove in \cref{thm:critical-points} in an abstract setting that
arclength constrained critical knots are actually also critical knots for the more
relaxed variational problem of fixed-length, but under the assumption of 
one more
order of differentiability of the critical knot. 
 For the concrete
energies treated above  we provide this additional regularity
by carefully analyzing the additional Lagrange parameter in the variational equation.
In fact, we prove full $C^\infty$-smoothness of critical knots
for all non-negative linear combinations
of bending energy $\BendingEnergy$ and tangent-point energies $\TP^{(p,2)}$ including
the cases when only one of these energies appears.
\begin{theorem}[Smoothness of critical points of total energies]
\label{thm:smoothness}
Let $\TL_{\nu\theta}\ceq\nu \, \BendingEnergy + \theta \, \TP^{(p,2)}$ for 
 $p\in (4,5)$, $\nu,\theta\ge 0$ with $\nu+\theta>0$,
and $\ML\in\{\AL^{1+s(\nu)},\AL^{1+s(\nu)}_\bz\}$ for a fixed point 
$\bz\in\R^n$, where $s(\nu)=1$ if $\nu >0$ and $s(\nu)=s(p)$ as in \cref{eq:s-TP} if
$\nu=0$. Then   
each critical knot $\g\in\ML$
of the restricted energy $\TL_{\nu\theta}|_{\ML}$ is of class $C^\infty(\R/\Z,\R^n)$.
\end{theorem}
If $ \nu>0 $ and $ \theta=0 $, then also $ \BendingEnergy $-critical immersions with possible self-intersections are smooth.\\
\Cref{thm:smoothness} also holds if one replaces the tangent-point energies
$\TP^{(p,2)}$ by the integral Menger curvature functionals $\intM^{(p,2)}$
for $p\in (\frac{7}{3},\frac{8}{3})$, or by O'Hara's energies
$E^{\alpha,1}$ for $\alpha\in (2,3)$; see \cite{freches_2025}. 
Critical points of these three energy families
are actually  also real-analytic, 
which follows from \Cref{thm:critical-points} and the regularity proofs
for the length-constrained critical points of  O'Hara's energies
in \cite{vorderobermeier_2020}, or of
integral Menger curvature in \cite{steenebruegge-vorderobermeier_2022}, the methods of which transfer to the tangent-point energies. 


\subsection{Outline of the paper}
In \Cref{sec:PS-abstract} we develop sufficient criteria for the  validity of the (PS)-condition
in the abstract setting of general submanifolds of Hilbert spaces, not only
for a single energy as in \Cref{thm:PS-abstract} but also for total energies with a leading order term and
some additional energy of lower order; see \Cref{cor:PS-total-energy}. Moreover, with
\Cref{thm:gradflow_subconvergence} we
provide a long-time existence result for the gradient flow and subconvergence
to critical points  for energies of class $C^{1,1}_\loc$. 
That this flow preserves symmetry is shown in 
\Cref{prop:symmetry-preserving-grad-flow}. Then we
deduce as simple
consequences the existence of minimizing (PS)-sequences and energy minimizers;
see \cref{koro:critpoint_lowerenergy,koro:minimizer_existence}, as well as \cref{cor:critical+minimizers}
 for an alternative proof 
that works for energies that are only $C^1$. \Cref{sec:arclength} contains the proofs
that the sets $\AL^{1+s}$ and $\AL^{1+s}_\bz$ defined in \cref{eq:ALs}
and \cref{eq:ALsz} are actually smooth submanifolds of the Hilbert space
$H^{1+s}(\R/\Z,\R^n)$ for any $s\in (\frac{1}{2},1]$ and
$\bz\in\R^n$ 
(\Cref{thm:arclength-manifold} and \Cref{thm:arclength-manifold-point}).
We also define suitable projections onto the respective tangent spaces that are needed
in the abstract \Cref{thm:PS-abstract}. Then we prove in \Cref{sec:AL-toplogy} that the open set
of regular knots strongly deformation retracts onto the arclength submanifolds (see \Cref{thm:strong-deformation-ALs}), 
and that connected components
of the arclength submanifolds correspond to 
different knot classes (\Cref{thm:components-knot-classes}). 
In \Cref{sec:AL-critical} we prove that
arclength constrained critical knots are also critical under the more
relaxed fixed-length constraint, still in the abstract general setting; see \Cref{thm:critical-points} and \Cref{cor:critical-points}. 


\Cref{sec:PS-bending+knot-energies} is dedicated to the proofs of
our main results on the Palais--Smale condition. We start in  \Cref{sec:PS-bending}, however,
 with the pure bending energy $\BendingEnergy$
as a prototype example, both for the verification of the (PS)-condition (\Cref{cor:PS_elastic}) and its consequences (arclength-constrained $\BendingEnergy$-critical 
points via minimization
or gradient flows, see \Cref{cor:elastic_crit_gradient-flow}). \Cref{sec:PS-bending} is meant as a
comprehensible blueprint for the more complicated proofs involving non-local terms
in the following \cref{sec:PS-elastic-knot} on total energies, and in
\cref{sec:tanpoint} on the pure tangent-point energies. 
\Cref{thm:bending-knot-energy,thm:bending-moebius-energy} are shown in \cref{sec:PS-elastic-knot}, whereas \Cref{thm:PS-TP} is proven in \cref{sec:tanpoint}.

Finally, \cref{sec:Smoothness} is devoted to the proof of \cref{thm:smoothness}. 
This regularity theorem follows from the combination of the smoothness result for critical points of the pure bending energy $ \BendingEnergy $ (\cref{thm:elastic-smoothness-new}) with those for critical knots of total energies where $ \BendingEnergy $ is the leading order term (\cref{thm:smoothness-T-critical-new}) and of the pure tangent-point energies $ \TPp $ (\cref{thm:smoothness-TP-new}). 

For notational convenience, we use combinations of the indices $\inj,\reg,\mathrm{a}$ to denote \underline{i}njective,
\underline{r}egular, or \underline{a}rclength parametrized
curves in a function space 
$\mathscr{F}$, e.g., $\FS_\inj,$ $\FS_\reg,$ $\FS_\injreg, $
or $\FS_\injunitspeed$.
Furthermore, we denote by $ L(V,W) $ the space of bounded linear maps between normed spaces, abbreviating $ L(V,V)=L(V) $. Unless otherwise specified, $ \left\langle \cdot,\cdot \right\rangle  $ and $ \left|\cdot \right| $, refer to the Euclidean inner product and its induced Euclidean norm, respectively.



